% optimization/alignment/00_context_and_problem-DE.tex
\section*{Motivation}
In praktischen Anwendungen werden Alignments nicht nur aus Daten rekonstruiert,
sondern auch entworfen, verändert und verbessert. Daraus entstehen
Optimierungsprobleme, in denen Alignment-Parameter an geometrische, physische
und regulatorische Bedingungen angepasst werden
\cite{Kufver2000RailwayAlignment,Klein1998Trassierung}.

\section{Allgemeines Optimierungsproblem}
\begin{definition}[Alignment-Optimierungsproblem]
Ein Alignment-Optimierungsproblem besteht aus einem parametrisierten
Alignment-Modell, einer Zielfunktion und einer Menge von Bedingungen.
\end{definition}
\begin{definition}[Zielfunktion]
Eine Zielfunktion ist eine Abbildung
\[
J:\mathcal A\to\R,
\]
die einem Alignment \(\mathcal A\) ein Kosten- oder Qualitätsmaß zuordnet.
\end{definition}
Typische Ziele sind die Minimierung von Krümmungsänderung, Abweichung von
Messdaten und Baukosten sowie die Maximierung des Komforts. Diese Zielklassen
entsprechen allgemeiner numerischer Optimierung
\cite{Nocedal2006NumericalOptimization} und eisenbahnspezifischen Ansätzen
\cite{Kufver2000RailwayAlignment}.

\section{Parametrisierung}
Die Parameterwahl ist für die Formulierung zentral.
\begin{definition}[Parametervektor]
Ein Alignment wird durch einen Parametervektor
\[
\mathbf x\in\R^n
\]
beschrieben, der Größen wie Elementlängen, Krümmungswerte,
Übergangsparameter und geometrische Bedingungen kodiert.
\end{definition}
Das dünn besetzte Alignment-Modell bietet eine natürliche elementweise
Parametrisierung.
\TODO{Ausdrückliche Parametersätze für feste und Übergangselemente.}

\section{Bedingungen}
\begin{definition}[Bedingung]
Eine Bedingung hat die Form
\[
g_i(\mathbf x)\leq0,\qquad h_j(\mathbf x)=0.
\]
\end{definition}
Bedingungen entstehen aus geometrischer Stetigkeit,
Engineering-Entwurfsregeln, physischen Grenzen und regulatorischen
Anforderungen.

\subsection{Geometrische Bedingungen}
\begin{itemize}
\item Stetigkeit von Position und Richtung,
\item Stetigkeit der Krümmung.
\end{itemize}

\subsection{Engineering-Bedingungen}
\begin{itemize}
\item Mindestradius,
\item maximale Überhöhung,
\item Grenzen der Krümmungsableitungen.
\end{itemize}
Diese Bedingungen spiegeln übliche Eisenbahnentwurfspraxis und Regeln aus
Literatur und Normen wider \cite{Klein1998Trassierung,EN13803_2017}.
\OPEN{Formale Darstellung von Engineering-Regeln als Bedingungen.}

\section{Bezug zur numerischen Rekonstruktion}
Die Bewertung der Zielfunktion erfordert die Rekonstruktion der
Alignment-Geometrie aus dem Parametervektor. Optimierung ist daher intrinsisch
mit numerischer Rekonstruktion verbunden und gehört numerisch in den weiteren
Kontext von Differentialgleichungen und Approximation
\cite{Hairer1993ODE,Press2007NumericalRecipes}.

\section{Sequentielle quadratische Programmierung}
Sequentielle quadratische Programmierung (SQP) ist ein leistungsfähiges
Verfahren für nichtlineare beschränkte Optimierungsprobleme
\cite{Nocedal2006NumericalOptimization}.
\begin{definition}[SQP-Iteration]
Ein SQP-Verfahren approximiert das ursprüngliche Problem durch eine Folge
quadratischer Teilprobleme.
\end{definition}
Jede Iteration löst eine quadratische Approximation der Zielfunktion unter
linearisierten Bedingungen.
\TODO{Ausdrückliche SQP-Formulierung für Alignment-Parameter.}
\RESEARCH{Anpassung von SQP an dünn besetzte Alignment-Strukturen.}

\section{Alignment-Optimierung als Steuerungsproblem}
Die Krümmungsfunktion \(\kappaf(s)\) kann als Steuerungsvariable eines
kontinuierlichen Systems interpretiert werden.
\begin{definition}[Steuerungsformulierung]
Das Alignment-Problem kann als Optimalsteuerungsproblem mit den
Zustandsvariablen \((\gamma,\theta)\) und der Steuerung \(\kappaf\) formuliert
werden.
\end{definition}
Diese Sicht ist mit Kufvers krümmungszentrierter Perspektive vereinbar
\cite{Kufver1997MathematicalDescription,Kufver2000RailwayAlignment}.
\RESEARCH{Verbindung zwischen Alignment-Optimierung und Optimalsteuerungstheorie.}

\section{Optimierung auf Netzebene}
Reale Eisenbahnsysteme umfassen nicht nur einzelne Alignments, sondern Netze
mit Verzweigungs- und Zusammenführungsstrukturen.
\RESEARCH{Erweiterung der Optimierung von einzelnen Alignments auf Netze.}
\OPEN{Behandlung von Weichen, Kreuzungen und parallelen Gleisen.}

\section{Robustheit und Unsicherheit}
Optimierung muss Unsicherheit der Eingabedaten berücksichtigen.
\RESEARCH{Robuste Optimierung unter unsicherer Geometrie.}
\OPEN{Abwägung zwischen Optimalität und Robustheit.}

\section{Implementierungssicht}
Praktische Optimierungssysteme müssen numerische Stabilität, Recheneffizienz
und Ausdrucksfähigkeit des Modells ausgleichen. Diese Abwägung ist in
allgemeiner numerischer \cite{Nocedal2006NumericalOptimization} und
eisenbahnspezifischer Optimierung \cite{Kufver2000RailwayAlignment} zentral.
\TODO{Integration der Optimierung in den Modellierungsablauf.}

\section{Historischer Kontext}
Klassische Systeme wie AXTRAN haben die Machbarkeit optimierungsbasierten
Alignment-Entwurfs gezeigt. Moderne Ansätze erweitern diese Ideen mit
flexibleren Modellen und verbesserten numerischen Verfahren im Übergang von
klassischem Eisenbahnentwurf zu ausdrücklicher Optimierungsformulierung
\cite{Kufver2000RailwayAlignment}.
\TODO{Rekonstruktion und Verallgemeinerung AXTRAN-artiger Verfahren.}
